% generated by Paperino/scripts/geometry_control.py -- do not edit by hand
\begin{table}[H]
\caption{Detection against a generator whose pointer moves take random steps.}
\label{tab:geometry}
\centering
\small
\setlength{\tabcolsep}{3pt}
\begin{adjustwidth}{-\extralength}{0cm}\resizebox{\linewidth}{!}{\begin{tabular}{lrrrrrrrr}
\toprule
 & & \texttt{step\_cv} & Kind-stratified & \multicolumn{2}{c}{Forest LOKO mean $F_1$} & \multicolumn{2}{c}{Rule LOKO mean $F_1$} & LOKO \\
\cmidrule(lr){5-6} \cmidrule(lr){7-8}
Automation class & Windows & $= 0$ (\%) & $F_1$ (95\% interval) & 0.5 & 1\% budget & \texttt{step\_cv} $= 0$ & No movement & AUC \\
\midrule
Released generator (fixed steps) & 645 & 97.8 & 0.988 (0.979--0.996) & 0.459 & 0.782 & 0.737 & 0.379 & 0.969--1.000 \\
Random steps (1 to 10\,px per axis) & 645 & 54.9 & 0.985 (0.974--0.993) & 0.439 & 0.742 & 0.451 & 0.379 & 0.967--0.999 \\
\bottomrule
\end{tabular}}\end{adjustwidth}
\par\vspace{2pt}\noindent\footnotesize Genuine windows are the released ones; only the automation class changes. The random-step generator keeps every event time, kind, key, and click of the released one and draws every pointer move as a random step from a separate seeded generator. Designs as in Tables~\ref{tab:stratified} and~\ref{tab:transfer}. The rules are those of Table~\ref{tab:baselines}, which learn nothing: \texttt{step\_cv} $= 0$ flags every window with zero step variation, and the no-movement rule every window without measurable pointer movement.
\end{table}
